\chapter{A Concrete Naming Certificate for $\RegularCauchyModulusIndependenceUp$}
\label{ch:concrete-instances-regular-cauchy-modulus-independence-namecert}
\origin{human}

Following Bishop and Bridges constructive analysis, the $\RegularCauchyModulusIndependenceUp$ packet records the finite BEDC route by which a single regular Cauchy rational source presents two displayed modulus ledgers while producing one $\RealUp$ seal. The chapter sits after the finite window discipline of \autoref{ch:concrete-instances-streamname-namecert}, the regular rational handoff of \autoref{ch:concrete-instances-regseqrat-namecert}, the dyadic tolerance ledger of \autoref{ch:concrete-instances-dyadicratcore-namecert}, the modulus realizer surface of \autoref{ch:concrete-instances-regular-cauchy-modulus-realizer-namecert}, the stability route of \autoref{ch:concrete-instances-regular-cauchy-modulus-stability-namecert}, and the terminal seal of \autoref{ch:concrete-instances-real-namecert}. Its purpose is to make the $\RealUp$ seal depend on the shared regular Cauchy source and the displayed cofinal comparison, not on an accidental choice of modulus ledger.

\begin{definition}[Regular Cauchy modulus-independence carrier]
\label{def:regular-cauchy-modulus-independence-carrier}
A $\RegularCauchyModulusIndependenceUp$ carrier is a finite $\BHist$ packet
$$
I=(R,S_0,S_1,D_0,D_1,J,A,H,C,P,N)
$$
with the following displayed rows:
\begin{enumerate}
\item $R$ is the single $\RegSeqRatUp$ source row, exposing the regular rational approximants that both modulus ledgers read.
\item $S_0$ and $S_1$ are the two $\StreamNameUp$ schedule rows used to request finite windows from the same source row $R$.
\item $D_0$ and $D_1$ are the two $\DyadicRatCoreUp$ modulus-ledger rows, including their displayed dyadic radii, slack comparisons, and finite precision requests.
\item $J$ is the cofinal comparison row, recording the finite zig-zag by which requests through $D_0$ and requests through $D_1$ are compared over the shared source row $R$.
\item $A$ is the unique $\RealUp$ seal row reached after the two schedules, two ledgers, and the cofinal comparison have been displayed.
\item $H$ records componentwise $\hsame$ transport for the source row, both schedule rows, both dyadic modulus ledgers, the cofinal comparison row, the Real seal, replay, provenance, and local naming rows.
\item $C$ records displayed $\Cont$ replay of the route
$$
R\longmapsto S_0,S_1\longmapsto D_0,D_1\longmapsto J\longmapsto A .
$$
\item $P$ is the $\Pkg$ provenance over $\RegularCauchyModulusIndependenceUp$, $\RegSeqRatUp$, $\StreamNameUp$, $\DyadicRatCoreUp$, $\RegularCauchyModulusRealizerUp$, $\RegularCauchyModulusStabilityUp$, $\RealUp$, $\BHist$, $\hsame$, $\Cont$, and $\NameCert$.
\item $N$ is the local $\NameCert_{\RegularCauchyModulusIndependenceUp}$ row.
\end{enumerate}
The classifier compares accepted packets only through $R,S_0,S_1,D_0,D_1,J,A,H,C,P,N$. It has no coordinate for a preferred modulus ledger, quotient stream equality, selected limits, countable choice, ambient $\RealUp$ completeness, Lean verification, or bridge checking.
\end{definition}

\begin{theorem}[Regular Cauchy modulus-independence NameCert obligations]
\label{thm:regular-cauchy-modulus-independence-namecert-obligations}
For every accepted $\RegularCauchyModulusIndependenceUp$ carrier
$$
I=(R,S_0,S_1,D_0,D_1,J,A,H,C,P,N),
$$
the local $\NameCert_{\RegularCauchyModulusIndependenceUp}$ surface consists exactly of the following finite obligation rows:
\begin{enumerate}
\item carrier admission for the shared $\RegSeqRatUp$ source row, the two $\StreamNameUp$ schedules, the two $\DyadicRatCoreUp$ modulus ledgers, the cofinal comparison row, the unique $\RealUp$ seal, and the structural rows $H,C,P,N$;
\item shared-source exposure, requiring both schedules and both ledgers to project the same source row $R$ before any seal-facing row is inspected;
\item two-modulus cofinal comparison, requiring $J$ to display finite comparison data from the $D_0$ route to the $D_1$ route and back over the source row $R$;
\item $\RegSeqRatUp$ handoff, requiring every ledger request to be read as a regular rational approximation from $R$ rather than from a detached stream representative;
\item $\RealUp$ seal factorization, requiring every seal consumer to pass through the shared source row and the cofinal comparison before reaching $A$;
\item transport, replay, and provenance, requiring $H,C,P,N$ to preserve, replay, source, and name only the displayed rows of the packet;
\item ledger non-escape, requiring the certificate to export no preferred modulus ledger, quotient stream equality, selected limit, countable-choice schedule, ambient $\RealUp$ completeness, Lean verification, or bridge checking.
\end{enumerate}
No modulus-independence obligation is stored outside these rows.
\end{theorem}
\begin{proof}
Project the accepted packet $I=(R,S_0,S_1,D_0,D_1,J,A,H,C,P,N)$. Carrier admission exposes one regular rational source row, two finite schedules, two dyadic modulus ledgers, one cofinal comparison row, one terminal Real seal, and the four structural rows.

The shared-source clause is read first. Both schedules $S_0,S_1$ request windows from $R$, and both dyadic ledgers $D_0,D_1$ are attached to those requests, so neither ledger can present an independent stream representative.

The cofinal clause is the row $J$. It records finite comparison data in both directions between the requests displayed by $D_0$ and $D_1$, always over the same source row $R$.

The handoff to $\RealUp$ is downstream. A seal-facing consumer must follow the route through $R$, the two schedules, the two ledgers, and $J$ before it can inspect $A$.

Finally, $H,C,P,N$ transport, replay, source, and name this route without adding an excluded coordinate. Since the carrier has no row for a preferred ledger or for an ambient completeness principle, the listed rows exhaust the local NameCert surface.
\end{proof}

\begin{theorem}[Regular Cauchy modulus-independence Real seal factorization]
\label{thm:regular-cauchy-modulus-independence-real-seal-factorization}
For every accepted $\RegularCauchyModulusIndependenceUp$ carrier
$$
I=(R,S_0,S_1,D_0,D_1,J,A,H,C,P,N),
$$
every $\RealUp$ consumer factors through
$$
R\longmapsto S_0,S_1\longmapsto D_0,D_1\longmapsto J\longmapsto A
$$
together with the structural rows $H,C,P,N$. In particular, replacing the selected finite request route from $D_0$ by the displayed cofinal route through $D_1$, or conversely, leaves the exposed $\RealUp$ seal row $A$ named by the same local certificate.
\end{theorem}
\begin{proof}
Start with a consumer of the terminal Real seal. The carrier projection exposes $A$, but the displayed route places $A$ after the two schedule rows, the two dyadic modulus ledgers, and the cofinal comparison row.

Trace the consumer backward from $A$. The only source coordinate available to either modulus ledger is $R$, and the only schedules available to the ledgers are $S_0$ and $S_1$.

The comparison step is $J$. A request that arrives through $D_0$ is compared to a request through $D_1$, and the converse comparison is also displayed, so the consumer receives a finite cofinal witness before the Real seal is read.

The structural rows preserve this factorization. The $\hsame$ row transports corresponding displayed components, the $\Cont$ row replays the same ordering, and the $\Pkg$ and $\NameCert$ rows name no extra modulus choice.

Thus the Real seal is a downstream endpoint of the shared regular-Cauchy source and its cofinal ledger comparison. It is not a certificate of a preferred modulus, quotient stream equality, selected limits, or ambient completeness.
\end{proof}

\closureat{RegularCauchyModulusIndependenceUp}{seedStr}

\begin{closurestatus}{\RegularCauchyModulusIndependenceUp}
  \constructivestory{A $\RegularCauchyModulusIndependenceUp$ packet is generated from one $\RegSeqRatUp$ source row, two $\StreamNameUp$ schedules, two $\DyadicRatCoreUp$ modulus ledgers, one cofinal comparison row, one $\RealUp$ seal, componentwise $\hsame$ transport, displayed $\Cont$ replay, $\Pkg$ provenance, and a local $\NameCert$ row.}
  \theoryclosure{\seedClosure}
  \scopeclosed{\autoref{def:regular-cauchy-modulus-independence-carrier}, \autoref{thm:regular-cauchy-modulus-independence-namecert-obligations}, and \autoref{thm:regular-cauchy-modulus-independence-real-seal-factorization} seed the finite L10 route separating Real seals from accidental modulus-ledger choices.}
  \formalstatus{\unformalizedV}
  \bridgestatus{none}
  \notclaimed{This chapter does not close preferred modulus selection, quotient stream equality, selected limits, countable choice, ambient $\RealUp$ completeness, Lean verification, or bridge checking.}
  \upgradepath{Obligation closure requires checked carrier admission, shared-source exposure, two-modulus cofinal comparison, RegSeqRat handoff, Real-seal factorization, transport and replay coverage, provenance coverage, and ledger non-escape for BEDC.Derived.RegularCauchyModulusIndependenceUp.}
\end{closurestatus}
